% Self-contained research section for integration into the parent article.
% Requires amsmath, amssymb, amsthm; standard theorem environments.
% Source target: ProveIt_Canonical_Digamma_Fifth_Moment_2026-10-11,
% sections/09_questions.tex, Conjecture conj:firstq; its inherited matrix
% normalization is sections/07_tail_rank.tex, equations tr:Aentries,
% tr:column and tr:operator. No source files have been modified.

\section{Optimal initial moments for adjacent even tail orders}
\label{ar:section}

The incoming canonical-digamma report \cite{Fifth} conjectures that the first $q$
centered moments of two unequal even harmonic tails have full formal
rank $q$ (its Conjecture~\texttt{conj:firstq}). That report proves the
claim when the two orders coincide and gives a larger sufficient prefix
for every unequal pair. We prove the conjectured optimal prefix for the
entire next family
\[
                 (p,q)=(2b-2,2b),\qquad b\ge2.
\]
The proof also gives an exact determinant with a strictly positive
rational recurrence. Thus the nonvanishing is an all-order theorem,
independent of numerical experiments. The conjecture for general larger
gaps remains open here.

\subsection{The inherited coordinates and the exact target}

For even integers $2\le p\le q=2b$, put
\[
 A_p(n)=\zeta(p)-H_n^{(p)},\qquad
 D_{p,q}(s)=\sum_{n\ge0}
     A_p(n)A_q(n)(n+\tfrac12)^{-s},\qquad
 D_r=D_{p,q}(-2r).
\]
The series initially defines $D_{p,q}$ for
$\Re s>3-p-q$; the notation $D_r$ uses its meromorphic continuation.
The centered asymptotic series for the product of the two tails contains
only even powers, so every point $s=-2r$ is regular. This regularity and
the following finite coordinate formulas are inherited from the two-tail
calculus, not additional conclusions of the rank argument.

The source supplies blocks $C_j$ with mutually inverse transforms
\begin{align}
 D_r&=\sum_{j=0}^r
    \frac{\binom{2r}{2j}}{2j+1}B_{2r-2j}(\tfrac12)C_j,
       \label{ar:forward}\\
 C_j&=\sum_{r=0}^j
       \binom{2j+1}{2r}2^{-2(j-r)}D_r.
       \label{ar:backward}
\end{align}
Here $B_n(x)$ is the Bernoulli polynomial and $B_n=B_n(0)$, with
$B_1=-1/2$. The first $b$ blocks $C_0,\ldots,C_{b-1}$ are kept as
specified coordinate labels. All subsequent blocks have the form
\begin{equation}\label{ar:block}
 C_j=\sum_{k=0}^{b-1}a_{jk}Z_k,\qquad
 Z_k=\zeta(2k),\quad Z_0=-\tfrac12,\qquad j\ge b.
\end{equation}
For $d=2j+1$ and $e=d-p-q+2k+1$, the exact rational coefficients are
\begin{equation}\label{ar:entries}
 a_{jk}=\begin{cases}
 B_e\left\{
 \displaystyle\frac{\binom{d-q+1}{p-2k}}{d-q+1}
 +\displaystyle\frac{\binom{d-p+1}{q-2k}}{d-p+1}
 \right\},&e\ge0,\\[6pt]
 0,&e<0.
 \end{cases}
\end{equation}
A binomial summand is omitted unless its lower index is at least two
and no larger than its upper index. In particular a putative lower
index zero is omitted. This endpoint convention is essential to the
proof and to the determinant.

We consider the $b\times b$ late-block matrix
\begin{equation}\label{ar:M}
 M_{p,q}=(a_{b+r-1,k})_{1\le r\le b,\ 0\le k<b}.
\end{equation}
The $2b\times2b$ coefficient matrix of
$D_0,\ldots,D_{2b-1}$ on the formal labels
\[
           C_0,\ldots,C_{b-1},Z_0,\ldots,Z_{b-1}
\]
is nonsingular exactly when $M_{p,q}$ is nonsingular. Indeed the
corresponding block matrix for $C_0,\ldots,C_{2b-1}$ is
$\operatorname{diag}(I_b,M_{p,q})$, and the transform
\eqref{ar:forward} is triangular with nonzero diagonal
$1/(2r+1)$. These are formal ranks over $\mathbb Q$; they imply
generation statements for the actual values without asserting their
arithmetic independence.

\subsection{A positive recurrence and the determinant}

Define positive rational numbers
\begin{equation}\label{ar:alpha}
 \alpha_n=\frac{2|B_{2n}|}{(2n)!}
         =\frac{4\zeta(2n)}{(2\pi)^{2n}},\qquad n\ge1.
\end{equation}
The equality follows from Euler's evaluation of even zeta values and
establishes positivity at every index. Define $\rho_1=1$ and, for $r\ge1$,
\begin{equation}\label{ar:rho}
 \rho_{r+1}=
 \frac{1}{4(2r+1)^2}
 \sum_{s=1}^r\alpha_{r+1-s}
  \left\{2r(2r+1)
       +\mathbf1_{s\ge2}\,2s(2s-1)\right\}\rho_s.
\end{equation}
Every coefficient in this recurrence is strictly positive. In particular
$\rho_s>0$ for every $s$. The beginning is
\[
 \rho_1=1,\qquad \rho_2=\frac1{36},\qquad
 \rho_3=\frac{11}{5400},\qquad
 \rho_4=\frac{113}{705600},\qquad
 \rho_5=\frac{2509}{190512000}.
\]

\begin{theorem}[The optimal prefix for adjacent even orders]
\label{ar:main}
For $q=2b\ge4$ and $p=q-2$, the matrix in \eqref{ar:M} has determinant
\begin{equation}\label{ar:det}
 \boxed{\displaystyle
 \det M_{q-2,q}
 =\frac{(-1)^{b(b-1)/2}}{b!}
       \bigl((2b+1)!!\bigr)^2\rho_{b+1}\ne0.}
\end{equation}
Consequently the first $q$ moments $D_0,\ldots,D_{q-1}$ have full
formal rank $q$, and
\begin{equation}\label{ar:spans}
 \operatorname{span}_{\mathbb Q}\{D_0,D_1,\ldots\}
 =\operatorname{span}_{\mathbb Q}\{D_0,\ldots,D_{q-1}\}
 =\operatorname{span}_{\mathbb Q}
     \{C_0,\ldots,C_{b-1},Z_0,\ldots,Z_{b-1}\}.
\end{equation}
The complete first-$q$ moment matrix, in the displayed coordinate order,
has determinant
\begin{equation}\label{ar:complete}
 \frac{(-1)^{b(b-1)/2}\bigl((2b+1)!!\bigr)^2\rho_{b+1}}
      {b!\,(4b-1)!!}.
\end{equation}
No prefix of fewer than $q$ rows can have rank $q$ on these $q$ formal
labels, so the prefix length is optimal in this specified formal system.
\end{theorem}

\subsection{Proof by rotation of the local differential equation}

We reproduce the source's differential description so that the argument
can be checked directly against \eqref{ar:entries}. Put $h(t)=\coth(t/2)$.
For $0\le k<b$, define the analytic germs
\begin{align}
 f_k(t)={}&\frac{1}{2(p-2k)!(q-1)!}
      \int_0^t u^{p-2k}(t-u)^{q-1}h(u)\,du\notag\\
 &+\frac{1}{2(q-2k)!(p-1)!}
      \int_0^t u^{q-2k}(t-u)^{p-1}h(u)\,du.
 \label{ar:column}
\end{align}
The first summand is omitted when $p-2k<2$. All included integrands
are regular at $u=0$. Their power series converge in $|t|<2\pi$.
Expanding
\[
        h(u)=2\sum_{n\ge0}\frac{B_{2n}}{(2n)!}u^{2n-1}
\]
and integrating monomials shows that the coefficient of
$t^{2j+1}/(2j+1)!$ in $f_k$ is precisely $a_{jk}$ for $j\ge b$.

For $v=(v_0,\ldots,v_{b-1})$ set
\[
 f_v=\sum_{k=0}^{b-1}v_kf_k,
 \qquad Q(t)=\sum_{k=0}^{b-1}v_k\frac{t^{q-2k}}{(q-2k)!}.
\]
When $p=q-2$, differentiating \eqref{ar:column} $q$ times gives
\begin{equation}\label{ar:operator}
 2f_v^{(q)}(t)
    =\bigl(Q''(t)-Q''(0)\bigr)h(t)+(Q(t)h(t))''
    =:\mathcal L_Q(t).
\end{equation}
The subtracted constant removes exactly the omitted lower-index-zero
term. Each side is an odd analytic germ at zero.

Rotate $t=ix$, put
\[
 R(x)=Q(ix)=\sum_{s=1}^b r_sx^{2s},\qquad
 r_s=\frac{(-1)^s}{(2s)!}v_{b-s},
 \qquad c(x)=\cot(x/2).
\]
Because $h(ix)=-ic(x)$ and two differentiations contribute a factor
$-1$, \eqref{ar:operator} gives
\begin{equation}\label{ar:rotated}
 \mathcal L_Q(ix)=i\bigl((R''-R''(0))c+(Rc)''\bigr).
\end{equation}
The benefit of this rotation is the one-sided coefficient expansion
\begin{equation}\label{ar:cot}
 c(x)=\frac2x-\sum_{n\ge1}\alpha_nx^{2n-1},\qquad |x|<2\pi,
\end{equation}
where every $\alpha_n$ is positive.

Let $Y_r$ be the ordinary coefficient of $x^{2r-1}$ in the bracket
on the right of \eqref{ar:rotated}. Write $r_{b+1}=0$ because $R$
has degree at most $2b$. Direct differentiation of the monomials in
\eqref{ar:cot} gives, for $1\le r\le b$,
\begin{equation}\label{ar:coefficient}
 Y_r=d_{r+1}r_{r+1}-\sum_{s=1}^r w_{rs}r_s,
\end{equation}
where
\begin{equation}\label{ar:weights}
 d_{r+1}=4(2r+1)^2,\qquad
 w_{rs}=\alpha_{r+1-s}
  \left\{2r(2r+1)+\mathbf1_{s\ge2}\,2s(2s-1)\right\}>0.
\end{equation}
For clarity, the coefficient supplied by $2/x$ and $r_sx^{2s}$,
$s\ge2$, is
\[
 2\{(2s)(2s-1)+(2s-1)(2s-2)\}r_s
       =4(2s-1)^2r_s.
\]
For $s=1$, the first differentiation term is removed by $R''(0)$,
while the second is the second derivative of a linear monomial and
vanishes. This explains both the superdiagonal $d_{r+1}$ and the
exceptional indicator in \eqref{ar:weights}.

Suppose first that $M_{q-2,q}v=0$. The first $b$ odd coefficients
of $\mathcal L_Q$ vanish, hence $Y_1=\cdots=Y_b=0$.
The first $b-1$ equations in \eqref{ar:coefficient} force
\[
                   r_s=\rho_sr_1,\qquad 1\le s\le b.
\]
The last equation, using $r_{b+1}=0$, then gives
\[
           0=-\sum_{s=1}^bw_{bs}\rho_sr_1
             =-d_{b+1}\rho_{b+1}r_1.
\]
All factors preceding $r_1$ are nonzero. Therefore $r_1=0$, then
$R=Q=0$ and $v=0$. This already proves nonsingularity over
$\mathbb C$, and therefore over $\mathbb Q$.

To retain the exact determinant, let $H_b$ be the matrix of
\eqref{ar:coefficient} from $(r_1,\ldots,r_b)$ to
$(Y_1,\ldots,Y_b)$. Replace its first column basis vector by
$(\rho_1,\ldots,\rho_b)^{\mathsf T}$, leaving the others fixed.
The basis change has determinant one. Its first column image has
only the final entry $-d_{b+1}\rho_{b+1}$. The first $b-1$ rows
on columns $2,\ldots,b$ form a lower triangular matrix with diagonal
$d_2,\ldots,d_b$. Expansion along the first column yields
\begin{equation}\label{ar:Hdet}
 \det H_b=(-1)^b\rho_{b+1}\prod_{s=2}^{b+1}d_s
    =(-1)^b4^b\bigl((2b+1)!!\bigr)^2\rho_{b+1}.
\end{equation}
On the other hand, coefficient comparison in \eqref{ar:rotated}
and \eqref{ar:operator} gives
\begin{equation}\label{ar:scalings}
 Y_r=\frac{2(-1)^{r-1}}{(2r-1)!}(M_{q-2,q}v)_r,
 \qquad r_s=\frac{(-1)^s}{(2s)!}v_{b-s}.
\end{equation}
The reversal of the $b$ columns has determinant
$(-1)^{b(b-1)/2}$. Combining the diagonal factors in
\eqref{ar:scalings} gives
\[
 \det M_{q-2,q}
     =\frac{(-1)^{b(b+1)/2}}{4^bb!}\det H_b.
\]
Substitute \eqref{ar:Hdet}; the resulting sign is
$(-1)^{b(b-1)/2}$, proving \eqref{ar:det}.
The block decomposition and triangular transform preceding the theorem
prove \eqref{ar:spans} and \eqref{ar:complete}.

\subsection{Explicit reconstruction without matrix inversion}

The same positive recurrence supplies an exact reconstruction algorithm.
Given $C_b,\ldots,C_{2b-1}$, define
\[
            Y_r=\frac{2(-1)^{r-1}}{(2r-1)!}C_{b+r-1},
                       \qquad 1\le r\le b.
\]
Put $\eta_1=0$ and compute, for $1\le r<b$,
\begin{equation}\label{ar:eta}
       \eta_{r+1}=\frac{Y_r+\sum_{s=1}^rw_{rs}\eta_s}{d_{r+1}}.
\end{equation}
Then
\begin{equation}\label{ar:reconstruct}
 \begin{split}
 r_1&=-\frac{Y_b+\sum_{s=1}^bw_{bs}\eta_s}
                  {d_{b+1}\rho_{b+1}},\\
 r_s&=\rho_sr_1+\eta_s\quad(1\le s\le b),\\
 Z_k&=(-1)^{b-k}(2b-2k)!\,r_{b-k}\quad(0\le k<b).
 \end{split}
\end{equation}
The equations follow by writing the solution of
\eqref{ar:coefficient} as $r_s=\rho_sr_1+\eta_s$, then using
the final row. The denominator in \eqref{ar:reconstruct} is strictly
positive by the theorem. Equations \eqref{ar:backward} and
\eqref{ar:reconstruct} express every $Z_k$ as an explicit rational
linear combination of the first $q$ moments. Substitution into
\eqref{ar:block} and \eqref{ar:forward} recovers every later moment.

For illustration, the first five late-block determinants are
\[
 \begin{array}{c|ccccc}
 (p,q)&(2,4)&(4,6)&(6,8)&(8,10)&(10,12)\\\hline
 \det M_{p,q}&-11/48&-113/384&2509/5120&309713/307200&-7097067/2867200.
 \end{array}
\]
Their nontrivial numerators now have an exact origin in the positive
recurrence \eqref{ar:rho}. Finite rational checks verify these examples,
the general determinant formula through the supplied test range, the
rotation factors, and the reconstruction algorithm. The infinite proof
is the coefficient argument above, not the finite check.

\subsection{An infinite family of ordinary rational evaluations}

The reconstruction has an arithmetic consequence at every adjacent pair,
not just at the example printed below. Let $q=2b$, $p=q-2$, and denote
the first row of $M_{p,q}^{-1}$ by $(\lambda_b,\ldots,\lambda_{q-1})$.
These rational coefficients are explicitly computed by
\eqref{ar:eta}--\eqref{ar:reconstruct}. Define
\begin{equation}\label{ar:all-weights}
 w_r=\sum_{j=\max\{b,r\}}^{q-1}
     \lambda_j\binom{2j+1}{2r}2^{-2(j-r)},\qquad
 P_b(x)=\sum_{r=0}^{q-1}w_rx^{2r}.
\end{equation}
Put
\[
 u_0=\frac1{(p-1)(q-1)},\qquad
 u_1=-\frac1{24}\left(\frac q{p-1}+\frac p{q-1}\right).
\]

\begin{corollary}[Rational evaluations from every adjacent pair]
\label{ar:all-rational}
For every even $q\ge4$, with $p=q-2$ and $x_n=n+1/2$,
\begin{equation}\label{ar:all-ordinary}
 \begin{split}
 \sum_{n=0}^{\infty}\bigg\{
 &\frac{P_b(x_n)}{(p-1)!(q-1)!}
       \psi^{(p-1)}(n+1)\psi^{(q-1)}(n+1)\\
 &-u_0w_{q-1}x_n^2
      -u_0w_{q-2}-u_1w_{q-1}\bigg\}=-\frac12.
 \end{split}
\end{equation}
The series is absolutely convergent. Its polynomial has exact degree
$2q-2$, since
\begin{equation}\label{ar:top-weight}
       w_{q-1}
       =\frac{b(4b-1)\rho_b}{(2b+1)^2\rho_{b+1}}>0.
\end{equation}
\end{corollary}

\begin{proof}
By the defining inverse row, $\sum_{j=b}^{q-1}\lambda_jC_j=Z_0=-1/2$.
The inverse Bernoulli transform \eqref{ar:backward} then gives
$\sum_{r=0}^{q-1}w_rD_r=-1/2$.
For centered $x=n+1/2$, the two leading terms of the tail product are
\[
                 A_p(n)A_q(n)=u_0x^{2-p-q}+u_1x^{-p-q}
                                  +O(x^{-p-q-2}).
\]
Because $p+q=2q-2$, multiplication by $P_b$ leaves precisely the
quadratic and constant terms subtracted in \eqref{ar:all-ordinary}.
The remainder is $O(x^{-2})$. Both subtracted terms have continued
centered sum zero, so the convergent series equals the specified moment
combination. The polygamma normalization follows from the two even
orders. Finally the coefficient of $C_{q-1}$ in the reconstruction
of $Z_0$ is $b\rho_b/((2b+1)^2\rho_{b+1})$, as follows directly
from the last row of \eqref{ar:reconstruct}. Multiplication by
$\binom{2q-1}{2q-2}=2q-1=4b-1$ proves \eqref{ar:top-weight}.
\end{proof}

\subsection{A convergent rational polygamma identity at orders four and six}

The first nontrivial new adjacent example after $(2,4)$ is $(4,6)$.
In this case \eqref{ar:entries} and the exact reconstruction give
\begin{equation}\label{ar:46matrix}
 \begin{pmatrix}C_3\\C_4\\C_5\end{pmatrix}
 =\begin{pmatrix}
 0&3/4&1/4\\5/12&2/3&-1/12\\1&-3/8&1/12
 \end{pmatrix}
 \begin{pmatrix}-1/2\\\zeta(2)\\\zeta(4)\end{pmatrix},
\end{equation}
and the inverse matrix is
\[
 \frac1{339}
 \begin{pmatrix}-28&180&264\\136&288&-120\\948&-864&360\end{pmatrix}.
\]
The first row, followed by \eqref{ar:backward}, yields
\begin{equation}\label{ar:46moment}
 67D_0+15516D_1+324320D_2+1434496D_3
       +1601280D_4+371712D_5=-21696.
\end{equation}
This has an ordinary, absolutely convergent realization. Put
\[
 \begin{split}
 P(x)={}&67+15516x^2+324320x^4+1434496x^6\\
        &+1601280x^8+371712x^{10},\qquad x_n=n+\tfrac12.
 \end{split}
\]
Then
\begin{equation}\label{ar:46ordinary}
 \boxed{\displaystyle
 \sum_{n=0}^{\infty}\left\{
  \frac{P(x_n)}{720}\psi^{(3)}(n+1)\psi^{(5)}(n+1)
       -\frac{123904}{5}x_n^2-\frac{316928}{5}\right\}=-21696.}
\end{equation}
To check the subtraction and the meaning of this identity, the centered
Hurwitz expansion gives
\[
 A_4(n)A_6(n)
   =\frac1{15}x_n^{-8}-\frac7{60}x_n^{-10}+O(x_n^{-12}).
\]
Multiplication by $P$ leaves precisely the two nonnegative powers
displayed in \eqref{ar:46ordinary}; the remaining bracket is
$O(x_n^{-2})$. Their continued centered sums are zero because
$\zeta(-2,1/2)=\zeta(0,1/2)=0$. Thus the absolutely convergent
sum equals the meromorphically continued moment combination in
\eqref{ar:46moment}. Finally
$A_4(n)=\psi^{(3)}(n+1)/6$ and
$A_6(n)=\psi^{(5)}(n+1)/120$, explaining the factor $720$.
Every summand is an ordinary real number; no finite-part instruction
remains in \eqref{ar:46ordinary}.

There is also a finite telescoping proof of \eqref{ar:46ordinary}
which is independent of the continued-moment formulas. Define
\begin{align*}
 K(N)&=-3584N^7+23040N^9+33792N^{11},\\
 P_4(N)&=128N(N+1)
       (44N^4+88N^3+67N^2+23N-14),\\
 P_6(N)&=128N^3(N+1)^3(33N^2+33N-14),\\
 R_0(N)&=-\frac{74624}{15}N^3+4928N^2-\frac{244288}{15}N.
\end{align*}
Put
\[
 \mathcal F(N)=K(N)A_4(N)A_6(N)
      +P_4(N)A_4(N)+P_6(N)A_6(N)+R_0(N).
\]
The identities $A_p(N-1)=A_p(N)+N^{-p}$, together with direct
polynomial multiplication, give
\begin{equation}\label{ar:46telescoping}
 \mathcal F(N)-\mathcal F(N-1)
  =P(N-\tfrac12)A_4(N-1)A_6(N-1)
      -\frac{123904}{5}(N-\tfrac12)^2-\frac{316928}{5}.
\end{equation}
This equality is a rational-function identity in $N$ and in two
indeterminates replacing $A_4(N)$ and $A_6(N)$; no zeta relation is
used to verify it. Since $\mathcal F(0)=0$, the partial sum of
\eqref{ar:46ordinary} through $n=N-1$ is exactly $\mathcal F(N)$.
The standard uncentered Hurwitz expansions
\[
 \begin{split}
 A_4(N)&=\frac1{3N^3}-\frac1{2N^4}+\frac1{3N^5}+O(N^{-7}),\\
 A_6(N)&=\frac1{5N^5}-\frac1{2N^6}+\frac1{2N^7}+O(N^{-9})
 \end{split}
\]
give $\mathcal F(N)=-21696+O(N^{-1})$. All positive powers
cancel; the constant contributions from $KA_4A_6$, $P_4A_4$, and
$P_6A_6$ are respectively $-20224$, $-448$, and $-1024$.
This proves the evaluation by an independent finite identity followed
by an elementary asymptotic limit.

\paragraph{Scope.}
This proves the source conjecture for two infinite classes: the source's
diagonal class and the new adjacent-order class. The phrase ``optimal
prefix'' refers to the formal coordinate presentation; it does not
prove that $q$ actual numerical constants are independent. For larger
gaps the rotated equation has several free initial coefficients.
Positivity of its scalar recurrence alone does not establish
nonsingularity of the resulting multi-column boundary matrix. That
remaining distinction is essential and prevents an unsupported
extension of the present theorem to all unequal pairs.

% Optional continuation after adjacent_tail_prefix.tex. All labels ar:.
\subsection{A smaller boundary determinant for a fixed gap}
\label{ar:general-gap}

The rotation argument has a useful exact extension even though it does
not settle the complete unequal-order conjecture. Write
$q=2b$, $p=2(b-c)$, with $1\le c<b$, so the differential order in
the source operator is $m=2c$. For integers $s\ge c+1$ put
\[
 A_s=\frac{(2s)!}{(2s-2c)!},\qquad
 B_s=\frac{(2s-1)!}{(2s-2c-1)!},\qquad
 d_s=2(A_s+B_s),
\]
and put $A_s=0$ for $1\le s\le c$. Define $c$ fundamental rational
sequences by
\begin{align}
 R_{s,j}&=\delta_{sj},&&1\le s,j\le c,\notag\\
 R_{s,j}&=\frac1{d_s}\sum_{u=1}^{s-1}
       \alpha_{s-u}(B_s+A_u)R_{u,j},
       &&s\ge c+1,\quad 1\le j\le c.
       \label{ar:fundamental}
\end{align}
Every generated entry $R_{s,j}$ is strictly positive. The relevant
boundary matrix has size $c$, independent of $b$:
\[
             F_{b,c}=(R_{b+i,j})_{1\le i,j\le c}.
\]

\begin{proposition}[Fixed-gap boundary reduction]\label{ar:gap-det}
In the normalization \eqref{ar:M},
\begin{equation}\label{ar:gap-factor}
 \det M_{2(b-c),2b}
  =\frac{(-1)^{b(b-1)/2}}{4^bb!}
       \left(\prod_{s=c+1}^{b+c}d_s\right)\det F_{b,c}.
\end{equation}
In particular the optimal-prefix conjecture for a given gap $2c$
is equivalent to $\det F_{b,c}\ne0$ for every integer $b>c$.
When $c=1$, $R_{s,1}=\rho_s$ and this is
Theorem~\ref{ar:main}.
\end{proposition}

\begin{proof}
For $R(x)=Q(ix)=\sum r_sx^{2s}$, rotate the source operator to
\[
 \bigl(R^{(2c)}-R^{(2c)}(0)\bigr)\cot(x/2)
        +(R\cot(x/2))^{(2c)}.
\]
The coefficient of $x^{2k-1}$, with $s=k+c$, is
\begin{equation}\label{ar:general-coefficient}
       d_sr_s-\sum_{u=1}^{s-1}\alpha_{s-u}(B_s+A_u)r_u.
\end{equation}
The pole $2/x$ gives $d_s$; every other cotangent coefficient has
the negative sign in \eqref{ar:cot}. The constant derivative
subtraction again makes $A_c=0$, and all lower derivatives already
vanish. This proves the displayed coefficient formula.

The first $b-c$ vanishing rows successively determine
$r_{c+1},\ldots,r_b$ from $r_1,\ldots,r_c$ by
\eqref{ar:fundamental}. The remaining $c$ rows, with
$r_{b+1}=\cdots=r_{b+c}=0$, are equivalent to the vanishing of
the boundary coordinates $F_{b,c}(r_1,\ldots,r_c)^{\mathsf T}$.
For the determinant, the first elimination contributes
$\prod_{s=c+1}^bd_s$. The remaining boundary block differs from
$-\operatorname{diag}(d_{b+1},\ldots,d_{b+c})F_{b,c}$ by a
unit lower triangular row transformation; its determinant is
$(-1)^c\prod_{s=b+1}^{b+c}d_s\det F_{b,c}$.
Moving the $c$ initial columns past the other $b-c$ columns contributes
$(-1)^{c(b-c)}$.

Finally, the rotation gives the row factor
$2(-1)^{r-c}/(2r-1)!$ and the same reversed column factors
$(-1)^s/(2s)!$ as before. The resulting conversion from the
coefficient matrix in \eqref{ar:general-coefficient} to $M$ has
factor $(-1)^{b(b-1)/2+bc}/(4^bb!)$. All additional signs multiply
to $(-1)^{c(1-c)}=1$, proving \eqref{ar:gap-factor}.
\end{proof}

The positivity of all individual entries in \eqref{ar:fundamental}
does not prove the determinant nonzero. There is not even a uniform
positive determinant sign for every fixed gap: exact evaluation at
$c=3$ gives $\det F_{6,3}<0$ and $\det F_{7,3}>0$. Both are nonzero.
The remaining research question is therefore a genuine boundary
determinant problem, rather than a direct extension of the one-dimensional
positivity proof.
